

%% 
%% This is a LaTeX document generated by automatic conversion of a
%% Mathematica notebook using Mathematica 4.0.
%% 
%% This document uses special macros that are defined in a LaTeX 
%% package file with a name of the form notebook*.sty.  Appropriate 
%% commands for loading the package have been included in this document.
%% The LaTeX package files can be found in a directory with path:
%% 
%% /usr/local/mathematica/SystemFiles/IncludeFiles/TeX/
%% 
%% The LaTeX package used in this document may require that your 
%% LaTeX implementation support fonts other than the standard Computer 
%% Modern fonts that are used by LaTeX.  See the Additional Information 
%% section under the TeXSave entry in the Mathematica Reference Guide 
%% for more information.
%% 
%% Implementation-specific instructions for installing TeX-related files
%% can be found at this URL.
%% 
%% http://www.wolfram.com/support/FrontEnds/Export/TeX
%% 
%% 

\documentclass{article}
\usepackage{notebook2e, latexsym}


\begin{document}

\Title{FOURIER TRANSFORM}
\Subtitle{LECTURE 13}
\Subsubtitle{ELE314: LINEAR SIGNALS AND SYSTEMS}
\Subsubtitle{Mariusz Jankowski \copyright{}

University of Southern Maine

mjkcc@usm.maine.edu}
\dispSFinmath{<<\Mvariable{Graphics`Graphics`}}
\begin{CellGroup}
\Section*{1 Introduction}


The Fourier transform (more recently the discrete-time version of the Fourier transform) is the most important system analysis tool available to
engineers and scientists. An in-depth understanding of frequency domain methods is a prerequisite to a discussion of techniques of great practical
importance in signal processing, communications and control.


\end{CellGroup}
\begin{CellGroup}
\Section*{2 Fourier Transform Equations}


The Fourier transform pair is defined by the following two equations:

\inlineTFinmath{x(t)=\frac{1}{2\pi }\multsp \int _{-\infty }^{\infty }X(\multsp \omega )\multsp {{\ExponentialE \multsp }^{j\multsp \omega \multsp
t\multsp }}\DifferentialD \omega }



\inlineTFinmath{X(\omega )=\multsp \int _{-\infty }^{\infty }x(\multsp t)\multsp {{\ExponentialE \multsp }^{-j\multsp \omega \multsp t\multsp }}\DifferentialD
t}



The function \inlineTFinmath{X(\omega )} is referred to as the {\itshape Fourier transform} of \inlineTFinmath{x(t)} and is also known as the signal's
{\itshape frequency spectrum}. In general, the function X(\(\omega \)) is complex, so we typically refer to the magnitude and phase spectra of a
signal. Equation (1), which maps a frequency domain signal into the time domain, is known as the synthesis equation (or inverse Fourier transform),
while Equation (2) is a mapping from the time domain to the frequency domain. It is known as the analysis equation or simply, the Fourier transform.

What kind of signals have Fourier transforms?

The so-called Dirichlet conditions [Oppenheim, 289-290], state that absolutely integrable signals (\inlineTFinmath{\int _{-\infty }^{\infty }\VerticalSeparator
x(t)\VerticalSeparator \DifferentialD t<\infty }), that are continuous or have a finite number of discontinuities, have Fourier transforms.




\end{CellGroup}
\begin{CellGroup}
\Section*{3 Fourier Analysis}


Fourier's important contribution to the signal analysis methods that now carry his name, was the development of a representation of aperiodic signals
as linear combinations of complex exponentials. The basic idea was to consider an aperiodic signal as a limiting case of a periodic signal with an
arbitrarily large period. (For a detailed discussion and derivation read Oppenheim/Willsky, pp.285-289.) As a result, the discrete line spectrum
of periodic signals, merges into a continuum of frequencies (the spacing between harmonically related sinusoids tends to zero).

Consider the Fourier transform of the following ideal pulse. 

\inlineTFinmath{x(t)\multsp =\multsp \big\{\MathBegin{MathArray}[c]{cc}
	1,&\big|t\big|\leq \frac{1}{2} \\
	0,&\Mvariable{else} 
  \MathEnd{MathArray}}

(Recall that a periodic pulse train consisting of such pulses was analyzed in \buttonbox[Hyperlink]{Lecture 9}: Fourier series.)



\begin{CellGroup}
\dispSFinmath{\MathBegin{MathArray}{l}
\Mfunction{Plot}[\Mfunction{UnitStep}[t+1/2]-\Mfunction{UnitStep}[t-1/2],  \\
\noalign{\vspace{0.5ex}}
\hspace{1.em} \{t,-2,2\},\Mvariable{PlotLabel}\rightarrow \Mvariable{Se\NTilde al\multsp aperiodica}]\\
\MathEnd{MathArray}}
\mathGraphic{lecture13.tex_gr1.eps}
\dispSFoutmath{\SkeletonIndicator \Mvariable{Graphics}\SkeletonIndicator }

\end{CellGroup}
Here we calculate the Fourier transform.By definition,the Fourier transform of this signal is:

\begin{CellGroup}
\dispSFinmath{\int _{-0.5}^{0.5}\exp [-\imag \multsp \omega \multsp t]\DifferentialD t\multsp //\Mfunction{ExpToTrig}}
\dispSFoutmath{\frac{2\multsp \sin [0.5\multsp \omega ]}{\omega }}

\end{CellGroup}
\dispSFinmath{\Muserfunction{X1}[\Mvariable{\omega \_}]:=\frac{2\sin \big[\multsp \frac{\omega }{2}\big]}{\multsp \omega }}
\begin{CellGroup}
\dispSFinmath{\MathBegin{MathArray}{l}
\Mfunction{Plot}[\Muserfunction{X1}[\omega ],\{\omega ,-15\multsp \pi ,15\multsp \pi \},\Mvariable{PlotRange}\rightarrow \Mvariable{All},  \\
\noalign{\vspace{0.5ex}}
\hspace{1.em} \Mvariable{PlotLabel}\rightarrow \Mvariable{Espectro\multsp frecuencial\multsp de\multsp un\multsp pulso}];\\
\MathEnd{MathArray}}
\mathGraphic{lecture13.tex_gr2.eps}

\end{CellGroup}
\begin{CellGroup}
\Exercise{Problem 1}
\ExerciseText{Find the zeros of the sinc function by hand calculation.}

\end{CellGroup}
Now consider a "smooth" function of time. Here we obtain the Fourier transform and later compare with the spectrum of the ideal square pulse. 

\dispSFinmath{\MathBegin{MathArray}{l}
\Muserfunction{x2}[\Mvariable{t\_}]:=0;  \\
\noalign{\vspace{0.9375ex}}  \\
\Muserfunction{x2}[\Mvariable{t\_}]:=\frac{1}{2}(\cos [\pi \multsp t]+1)\multsp /;\multsp \Mfunction{Abs}[t]<1
\MathEnd{MathArray}}
\begin{CellGroup}
\dispSFinmath{\Mfunction{Plot}[\Muserfunction{x2}[t],\{t,-4,4\},\Mvariable{PlotLabel}\rightarrow \Mvariable{Se\NTilde al\multsp suave}];}
\mathGraphic{lecture13.tex_gr3.eps}

\end{CellGroup}
\begin{CellGroup}
\dispSFinmath{\int _{-1}^{1}\frac{1}{2}(\cos [\pi \multsp t]+1)\multsp {{\ExponentialE }^{-\ImaginaryI \multsp \omega \multsp t}}\DifferentialD t//\Mfunction{ExpToTrig}\multsp
//\Mfunction{FullSimplify}}
\dispSFoutmath{\frac{{{\pi }^2}\multsp \sin [\omega ]}{{{\pi }^2}\multsp \omega -{{\omega }^3}}}

\end{CellGroup}
\dispSFinmath{\Muserfunction{X2}[\Mvariable{\omega \_}]:=\frac{{{\pi }^2}\multsp \sin [\omega ]}{{{\pi }^2}\multsp \omega -{{\omega }^3}}}
\begin{CellGroup}
\dispSFinmath{\Mfunction{Plot}[\Muserfunction{X2}[\omega ],\{\omega ,-15\multsp \pi ,15\multsp \pi \},\Mvariable{PlotRange}\rightarrow \Mvariable{All}];}
\mathGraphic{lecture13.tex_gr4.eps}

\end{CellGroup}

\end{CellGroup}
\begin{CellGroup}
\Section*{4 Visualization of Fourier Spectra}


Typical Fourier spectra are complex functions of frequency. There are a number of common techniques of plotting and visualizing the spectrum of a
signal. In contrast, the Fourier transform of (3) was real due to the even symmetry of signal \inlineTFinmath{x(t)}.  In this section, we will consider
the various methods of visualizing complex Fourier spectra. We begin by defining a well-known function of frequency, the so-called Butterworth filter
of order 3 (see Oppenheim/Willsky, p. 705). We first plot the real and imaginary components of this function. 



\begin{CellGroup}
\dispSFinmath{\MathBegin{MathArray}{l}
H[\Mvariable{\omega \_}]:=\frac{{{{{\omega }_c}}^3}}{{s^3}+2\multsp {{\omega }_c}\multsp {s^2}+2\multsp {{{{\omega }_c}}^2}\multsp s+{{{{\omega }_c}}^3}}/.\{s\rightarrow
\imag \multsp \omega ,{{\omega }_c}\rightarrow 2\multsp \pi \multsp 1000.\};  \\
\noalign{\vspace{0.90625ex}}  \\
H[\omega ]
\MathEnd{MathArray}}
\dispSFoutmath{\frac{2.4805\times {{10}^{11}}}{2.4805\times {{10}^{11}}+7.89568\times {{10}^7}\multsp \ImaginaryI \multsp \omega -12566.4\multsp
{{\omega }^2}-\ImaginaryI \multsp {{\omega }^3}}}

\end{CellGroup}
\begin{CellGroup}
\dispSFinmath{\MathBegin{MathArray}{l}
\Mvariable{DisplayTogetherArray}[  \\
\noalign{\vspace{0.5ex}}
\hspace{1.em} \{\Mfunction{Plot}[\Mfunction{Re}[H[\omega ]],\{\omega ,0,5000\multsp \pi \}],\Mfunction{Plot}[\Mfunction{Im}[H[\omega ]],\{\omega
,0,5000\multsp \pi \}]\}];\\
\MathEnd{MathArray}}
\mathGraphic{lecture13.tex_gr5.eps}

\end{CellGroup}
Frequently, the {\itshape magnitude} (sometimes, magnitude squared), the {\itshape phase} or both are of interest.



\begin{CellGroup}
\dispSFinmath{\MathBegin{MathArray}{l}
\Muserfunction{DisplayTogetherArray}[\{\Mfunction{Plot}[\Mfunction{Abs}[H[\omega ]],\{\omega ,0,5000\multsp \pi \},  \\
\noalign{\vspace{0.5ex}}
\hspace{3.em} \Mvariable{PlotLabel}\rightarrow \Mvariable{Magnitud\multsp espectral},\Mvariable{PlotRange}\rightarrow \Mvariable{All}],  \\
\noalign{\vspace{0.5ex}}
\hspace{2.em} \Mfunction{Plot}[\Mfunction{Arg}[H[\omega ]],\{\omega ,0,5000\multsp \pi \},\Mvariable{PlotLabel}\rightarrow \Mvariable{Fase\multsp
espectral}]\}];\\
\MathEnd{MathArray}}
\mathGraphic{lecture13.tex_gr6.eps}

\end{CellGroup}
Finally, another common method of displaying signal spectra is by using a {\itshape Bode plot}, which basically is a logarithmicaly scaled magnitude
squared plot. Here is the Bode plot of the example signal.



\begin{CellGroup}
\dispSFinmath{\Mfunction{Plot}[20\multsp \log [10,\Mfunction{Abs}[H[\omega ]]],\{\omega ,0.1,50000\}];}
\mathGraphic{lecture13.tex_gr7.eps}

\end{CellGroup}
Here is a classical Bode plot, with fancy grid lines, title, and logarithmicaly scaled frequency axis.

\begin{CellGroup}
\dispSFinmath{\MathBegin{MathArray}{l}
\Muserfunction{LogLinearPlot}[20\multsp \log [10,\Mfunction{Abs}[H[\omega ]]],  \\
\noalign{\vspace{0.5ex}}
\hspace{1.em} \{\omega ,0.1,100000\},\Mvariable{PlotStyle}\rightarrow \Mfunction{Thickness}[0.01],\multsp   \\
\noalign{\vspace{0.5ex}}
\hspace{1.em} \Mvariable{PlotLabel}->\Mvariable{Filtro\multsp Butterworth\multsp de\multsp orden\multsp 3},\Mvariable{GridLines}\rightarrow \Mvariable{Automatic}];\\
\MathEnd{MathArray}}
\mathGraphic{lecture13.tex_gr8.eps}

\end{CellGroup}

\end{CellGroup}
\begin{CellGroup}
\Section*{5 Fourier Synthesis}


Given the Fourier spectrum of a signal, we can use (1) to map the signal back into the time domain. This is known as the{\itshape  inverse Fourier
transform}. Consider the following simple, but IMPORTANT function of frequency (it is frequently called the {\itshape ideal lowpass filter;} the
meaning of this term will be explained shortly):



\dispSFinmath{X[\Mvariable{\omega \_}]:=\multsp \Mfunction{If}[\Mfunction{Abs}[\omega ]\multsp <\multsp 2\pi ,\multsp 1,\multsp 0]}
Following the definition of the inverse transform (2) we get:

\begin{CellGroup}
\dispSFinmath{\Mvariable{sig}=\frac{1}{2\pi }\int _{-2\pi }^{2\pi }{{\ExponentialE }^{\multsp \ImaginaryI \multsp \omega \multsp t}}\DifferentialD
\omega \multsp //\Mfunction{ExpToTrig}}
\dispSFoutmath{\frac{\sin [2\multsp \pi \multsp t]}{\pi \multsp t}}

\end{CellGroup}
\begin{CellGroup}
\dispSFinmath{\MathBegin{MathArray}{l}
\Muserfunction{DisplayTogetherArray}[\{\{\Mfunction{Plot}[X[\omega ],\{\omega ,-6\pi ,6\pi \}]\},  \\
\noalign{\vspace{0.5ex}}
\hspace{2.em} \{\Mfunction{Plot}[\Mvariable{sig},\{t,-5,5\},\Mvariable{PlotRange}\rightarrow \Mvariable{All}]\}\}];\\
\MathEnd{MathArray}}
\mathGraphic{lecture13.tex_gr9.eps}

\end{CellGroup}
Note the similarity between this Fourier pair, and the one in the section on Fourier analysis. The similarity is not coincidental and is an example
of a general property of the Fourier transform called duality (read Oppenheim/Willsky, pp. 208-211). This property is a direct result of the symmetry
in the definitions of the forward and inverse Fourier transforms.


\end{CellGroup}
\begin{CellGroup}
\Section*{6 Signal Bandwidth}


One very important concept in the theory of LTI systems is that of bandwidth. There is no single precise mathematical definition of bandwidth. In
practice, all methods use imprecise, qualitative notions of signal strength. The bandwidth of a signal may be defined as the range of frequencies
where the signal magnitudes are relatively "large". Consider the signal spectra \inlineTFinmath{\Muserfunction{X1}(\omega )} and \inlineTFinmath{\Muserfunction{X2}(\omega
)}. According to the description of bandwidth given above, which of the two given signals x1 and x2, has larger bandwidth? Here is a plot of the
spectra of the two signals (\inlineTFinmath{\Muserfunction{X1}(\omega )} is in red and \inlineTFinmath{\Muserfunction{X2}(\omega )} is in blue).




\begin{CellGroup}
\dispSFinmath{\MathBegin{MathArray}{l}
\Mvariable{DisplayTogetherArray}[  \\
\noalign{\vspace{0.5ex}}
\hspace{1.em} \{\{\Mfunction{Plot}[\Mfunction{Abs}[\Muserfunction{X1}[\omega ]],\multsp \{\omega ,\multsp -10\pi ,\multsp 10\pi \},\multsp \Mvariable{PlotRange}->\Mvariable{All}]\},
 \\
\noalign{\vspace{0.5ex}}
\hspace{2.em} \{\Mfunction{Plot}[\Mfunction{Abs}[\Muserfunction{X2}[\omega ]],\{\omega ,-10\pi ,10\pi \},\Mvariable{PlotRange}\rightarrow \Mvariable{All}]\}\}];\\
\MathEnd{MathArray}}
\mathGraphic{lecture13.tex_gr10.eps}

\end{CellGroup}
\Exercise{Problem 3}
\ExerciseText{Here we will compare the bandwidth of the two signals using the following metric. Define bandwidth as the range of frequencies that
include 98\%{} of the total signal energy. Begin by calculating the total signal energy, \inlineTFinmath{{e_T}}. Subsequently, by trial and error
estimate the "cutoff" frequency.}


*** UNFINISHED ***

time-bandwidth product

half-power point (use Butterworth filter as example)


\end{CellGroup}
\begin{CellGroup}
\Section*{7 Frequency Response of LTI Systems}


Linear time-invariant systems are frequently described by linear, constant-coefficient differential equations. The frequency response of such a system
may be obtained directly from the differential equation. Consider the simple case of the lowpass RC filter given by the differential equation,

\inlineTFinmath{a\multsp y'(t)+y(t)=x(t)}

Taking the Fourier transform of the expression given above, and simplifying, we obtain the frequency response of the system \inlineTFinmath{H(\omega
)}. Note particularly the use of the Fourier transform differentiation rule in going from (5) to (6).



\inlineTFinmath{\ScriptCapitalF [a\multsp y'(t)]+\ScriptCapitalF [y(t)]=\ScriptCapitalF [x(t)]}

\inlineTFinmath{a\multsp j\multsp \omega \multsp \multsp Y(\omega )\multsp +\multsp Y(\omega )\multsp =\multsp X(\omega )}

\inlineTFinmath{H(\omega )\multsp =\multsp \frac{Y(\omega )}{X(\omega )}\multsp =\multsp \frac{1}{1+j\multsp \omega \multsp a}}

Here we plot the magnitude and phase spectra of the first-order lowpass filter (we set \inlineTFinmath{a=1}).



\begin{CellGroup}
\dispSFinmath{\MathBegin{MathArray}{l}
\Muserfunction{DisplayTogetherArray}\big[\big\{\big\{\Mfunction{Plot}\big[\Mfunction{Abs}\big[\frac{1}{1+\imag \multsp \multsp \omega }\big],\{\omega
,0,5\pi \},\Mvariable{PlotRange}\rightarrow \Mvariable{All}\big]\big\},  \\
\noalign{\vspace{1.17708ex}}
\hspace{2.em} \big\{\Mfunction{Plot}\big[\Mfunction{Arg}\big[\frac{1}{1+\imag \multsp \multsp \omega }\big],\{\omega ,0,5\multsp \pi \},\Mvariable{PlotRange}\rightarrow
\Mvariable{All}\big]\big\}\big\}\big];\\
\MathEnd{MathArray}}
\mathGraphic{lecture13.tex_gr11.eps}

\end{CellGroup}
The frequency response of a system \inlineTFinmath{H(\omega )}, is the Fourier transform of its impulse response \inlineTFinmath{h(t)}, namely \inlineTFinmath{\multsp
h(t)\overvar{\rightarrow }{\ScriptCapitalF }H(\omega )}. Therefore, given the frequency response of a system, we can obtain the impulse response
using the inverse Fourier transform. (In version 3.x obtain the inverse transform by table look-up. The relevant transform pair is \inlineTFinmath{{e^{\multsp
a\multsp t\multsp }}u(t)\longrightarrow \multsp \frac{1}{a+j\multsp \omega }}, with \inlineTFinmath{\Mfunction{Re}[a]\multsp >\multsp 0}.)



\begin{CellGroup}
\dispSFinmath{\Muserfunction{InverseFourierTransform}\big[\frac{1}{1+\imag \multsp \multsp \omega },\omega ,t,\Mvariable{FourierParameters}\rightarrow
\{1,-1\}\big]}
\dispSFoutmath{{{\ExponentialE }^t}\multsp \Mfunction{UnitStep}[-t]}

\end{CellGroup}
\begin{CellGroup}
\Exercise{Problem 4}
\ExerciseText{Obtain the frequency response and the impulse response of an LTI system defined by the following differential equation: }
\ExerciseText{\inlineTFinmath{y''(t)+4\multsp y'(t)\multsp +\multsp 3\multsp y(t)\multsp =\multsp x'(t)\multsp +\multsp 2\multsp x(t)}}



\end{CellGroup}

\end{CellGroup}
\begin{CellGroup}
\Section*{PROBLEM SOLUTIONS}
\begin{CellGroup}
\Exercise{Problem 1}
\ExerciseText{Find the zeros of the sinc function by hand calculation.}

\end{CellGroup}
\begin{CellGroup}
\Exercise{Problem 2}
\ExerciseText{Consider the signal spectra \inlineTFinmath{\Muserfunction{X1}(\omega )} and \inlineTFinmath{\Muserfunction{X2}(\omega )}. According
to the description of bandwidth given above, which of the two given signals x1 and x2, has larger bandwidth? Here is a plot of the spectra of the
two signals (\inlineTFinmath{\Muserfunction{X1}(\omega )} is in red and \inlineTFinmath{\Muserfunction{X2}(\omega )} is in blue). }


\begin{CellGroup}
\dispSFinmath{\MathBegin{MathArray}{l}
\Mfunction{Plot}\big[\{\Muserfunction{X1}[\omega ],\Muserfunction{X2}[\omega ]\},\multsp \{\omega ,\multsp -10\pi ,\multsp 10\pi \},\multsp \Mvariable{PlotRange}->\Mvariable{All},
 \\
\noalign{\vspace{0.9375ex}}\Mvariable{PlotStyle}->\big\{\Mfunction{Hue}[0],\Mfunction{Hue}\big[\frac{2}{3}\big]\big\}\big];\\
\MathEnd{MathArray}}
\mathGraphic{lecture13.tex_gr12.eps}

\end{CellGroup}

\end{CellGroup}
\begin{CellGroup}
\Exercise{Problem 3}
\ExerciseText{Here we will compare the bandwidth of the two signals using the following metric. Define bandwidth as the range of frequencies that
include 98\%{} of the total signal energy. Begin by calculating the total signal energy, \inlineTFinmath{{e_T}}. Subsequently, by trial and error
estimate the "cutoff" frequency.}


Solution. First we evaluate total energy:

\begin{CellGroup}
\dispSFinmath{\big\{\Mvariable{E1}=1,\multsp \Mvariable{E2}=\frac{1}{4}\int _{-1}^{1}{{(\cos [\pi \multsp t]+1)}^2}\DifferentialD t\big\}}
\dispSFoutmath{\big\{1,\frac{3}{4}\big\}}

\end{CellGroup}
Next we evaluate \inlineTFinmath{{e_A}=\frac{1}{\pi }\int _{0}^{w}{{X(\omega )}^2}\DifferentialD \omega } for various values of the cutoff frequency
\inlineTFinmath{w} until \inlineTFinmath{{e_A}\multsp \approx \multsp 0.98\multsp {e_T}}.



\dispSFinmath{{e_A}[\Mvariable{X\_},\Mvariable{w\_}]:=\frac{1}{\pi }\Mfunction{NIntegrate}[{X^2},\multsp \{\omega ,0,\multsp w\}]}
In the case of signal x1 (i.e. unit pulse) we get

\begin{CellGroup}
\dispSFinmath{\Mfunction{Abs}\big[0.98-\Mfunction{Table}\big[{e_A}\big[\frac{2\sin \big[\multsp \frac{\omega }{2}\big]}{\multsp \omega },w\big],\{w,32,33,0.1\}\big]\big]}
\dispSFoutmath{\MathBegin{MathArray}{l}
\{0.000203173,0.000191107,0.000175227,  \\
\noalign{\vspace{0.5ex}}
\hspace{1.em} 0.000155127,0.000130451,0.000100897,0.0000662173,  \\
\noalign{\vspace{0.5ex}}
\hspace{1.em} 0.0000262237,0.0000192145,0.0000701685,0.00012665\}\\
\MathEnd{MathArray}}

\end{CellGroup}
\begin{CellGroup}
\dispSFinmath{\Mfunction{First}[\Mfunction{Position}[\%,\min [\%]]]*0.1+32}
\dispSFoutmath{\{32.9\}}

\end{CellGroup}
In the case of signal x2 (i.e. cosine shaped pulse) we get

\begin{CellGroup}
\dispSFinmath{\Mfunction{Abs}\big[0.98-\Mfunction{Table}\big[{e_A}\big[\frac{{{\pi }^2}\multsp \sin [\omega ]}{{{\pi }^2}\multsp \omega -{{\omega
}^3}},w\big]\Big/\Mvariable{E2},\{w,4,4.1,0.01\}\big]\big]}
\dispSFoutmath{\MathBegin{MathArray}{l}
\{0.00248435,0.00209329,0.00170761,  \\
\noalign{\vspace{0.5ex}}
\hspace{1.em} 0.00132728,0.000952236,0.000582436,0.000217827,  \\
\noalign{\vspace{0.5ex}}
\hspace{1.em} 0.000141641,0.000496018,0.000845353,0.0011897\}\\
\MathEnd{MathArray}}

\end{CellGroup}

\end{CellGroup}
\begin{CellGroup}
\dispSFinmath{\Mfunction{First}[\Mfunction{Position}[\%,\min [\%]]]*0.01+4.}
\dispSFoutmath{\{4.08\}}

\end{CellGroup}
\begin{CellGroup}
\Exercise{Problem 4}
\ExerciseText{Obtain the frequency response and the impulse response of an LTI system defined by the following differential equation: }
\ExerciseText{\inlineTFinmath{y''(t)+4\multsp y'(t)\multsp +\multsp 3\multsp y(t)\multsp =\multsp x'(t)\multsp +\multsp 2\multsp x(t)}}



\end{CellGroup}

\end{CellGroup}
\begin{CellGroup}
\Section*{HOMEWORK PROBLEMS}
\begin{CellGroup}
\Subsubsection*{Homework Problem 13.1}
Calculate frequency response of the following signal (see Oppenheim/Willsky, Example 4.1, pp. 290-291):

\dispSFinmath{\MathBegin{MathArray}{l}
\Mfunction{Clear}[s];  \\
\noalign{\vspace{0.9375ex}}  \\
s[\Mvariable{t\_}]:=\frac{1}{2}\multsp \exp \big[-\frac{t}{2}\big]\multsp \Mfunction{UnitStep}[t]
\MathEnd{MathArray}}

\end{CellGroup}
\begin{CellGroup}
\Subsubsection*{Homework Problem 13.2}
Here is a variant of the unit pulse in (3):

\inlineTFinmath{x(t)\multsp =\multsp \big\{\MathBegin{MathArray}[c]{cc}
	1,&0\leq t\multsp \leq 1 \\
	0,&\Mvariable{else} 
  \MathEnd{MathArray}}



 Calculate its Fourier transform. Compare the spectrum with that of signal (3). Discuss any differences and similarities.


\end{CellGroup}
\begin{CellGroup}
\Subsubsection*{Homework Problem 13.3}
(a) Calculate the Fourier transform of the following signal called the Gabor wavelet:

\begin{CellGroup}
\dispSFinmath{\MathBegin{MathArray}{l}
g[\Mvariable{t\_}]:=\exp [-{t^2}]\multsp \sin [2\multsp \pi \multsp a\multsp t]\multsp /.\multsp a->1.  \\
\noalign{\vspace{0.5ex}}  \\
\Mfunction{Plot}[g[t],\{t,-3,3\},\Mvariable{PlotLabel}\rightarrow \Mvariable{Gabor\multsp Wavelet}];
\MathEnd{MathArray}}
\mathGraphic{lecture13.tex_gr13.eps}

\end{CellGroup}
(b) Change (increase) the value of a and recalculate the Fourier transform. How does the Fourier transform change?


\end{CellGroup}
\begin{CellGroup}
\Subsubsection*{Homework Problem 13.4}
Calculate the bandwidth of the gabor wavelet. Consider bandwidth to be those frequencies, that are larger than 1/100 of the maximum magnitude value.


\end{CellGroup}

\end{CellGroup}

\end{document}
